\chapter{Models and Deployment}

\section{Role of IndicF5 and Vocos}

IndicF5 uses a pretrained Diffusion Transformer (DiT) with Continuous Flow
Matching. It predicts how the mel representation should change at a given
sampling time. The native sampler starts from Gaussian noise and follows these
predicted directions to generate acoustic features. The reference mel and
combined text provide the conditioning. Vocos decodes the resulting mel
features into complex frequency coefficients. The C++ engine reconstructs the
waveform from those coefficients with an inverse short-time Fourier transform.

\section{PyTorch export and INT8 quantization}

\texttt{export\_dit\_litert.py} wraps one IndicF5 velocity prediction at a
time and exports it as a fixed-shape FP32 FlatBuffer.
\texttt{export\_vocos\_litert.py} exports the Vocos decoder with separate
real and imaginary output planes. Exporting these two graphs separately lets
the native runtime control the sampling schedule and the inverse transform.

\texttt{quantize\_litert.py} applies the
\texttt{dynamic\_wi8c\_afp32} recipe. It stores eligible weights as
per-channel INT8 and retains floating-point audio and time interfaces and
integer text tokens. The packaged artifacts are
\texttt{indicf5\_dit\_int8.tflite}, \texttt{indicf5\_vocos\_int8.tflite}
and \texttt{vocab.txt}. The two model files occupy 341.20 MiB, compared with
1,339.54 MiB for the FP32 exports, a 74.5\% reduction on disk.

\clearpage
\section{Tensor interfaces and conversion validation}

\begin{table}[!ht]
\centering
\caption{Tensor interfaces for the fixed-window LiteRT graphs}
\label{tab:model-tensors}
\small
\begin{tabular}{@{}p{0.16\linewidth}p{0.43\linewidth}p{0.32\linewidth}@{}}
\toprule
Graph & Inputs & Output \\
\midrule
IndicF5 DiT &
Noisy mel and conditioning mel, both $[1,512,100]$ FP32. Text IDs
$[1,128]$ INT32 and time $[1]$ FP32. &
Velocity $[1,512,100]$ FP32. \\
Vocos & Mel $[1,100,512]$ FP32. &
Real and imaginary spectrum planes $[1,2,513,512]$ FP32. \\
\bottomrule
\end{tabular}
\end{table}

\subsection{Parity checks}

\texttt{evaluate\_parity.py} compares PyTorch and LiteRT FP32 using
synchronized inputs and sampling. The archived check reported 0.0037 dB of
dynamic-time-warping-aligned mel-cepstral distortion.
\texttt{benchmark\_int8.py} compares LiteRT INT8 against LiteRT FP32.
The archived eight-step check reported an average velocity cosine similarity
of 0.99923. These tests assess conversion fidelity. Listening quality,
intelligibility and speaker similarity require separate evaluations.
